\begin{summarybox}
	\textbf{\textcolor{CSummary}{一句话核心}}
	\textbf{向量数乘的几何意义可概括为：共线、模长按绝对值比例变化，方向由数乘系数的符号决定。}

	\medskip
	\textbf{\textcolor{CSummary}{使用条件}}
	\begin{itemize}[leftmargin=*, itemsep=0.5em, topsep=0.4em]
		\item \textbf{条件1（非零向量）：} 讨论方向关系和唯一表示时，默认 $\vec{a}\neq\vec{0}$。
		\item \textbf{条件2（实数任意）：} $\lambda\in\mathbb{R}$，可以取正数、负数或 $0$。
	\end{itemize}

	\medskip
	\textbf{\textcolor{CSummary}{关键提醒}}
	\begin{itemize}[leftmargin=*, itemsep=0.5em, topsep=0.4em]
		\item \textbf{易错点（漏绝对值）：} 计算 $|\lambda\vec{a}|$ 时注意 $|\lambda\vec{a}| = |\lambda|\,|\vec{a}|$，别写成 $\lambda|\vec{a}|$。
		\item \textbf{检查项（零向量特例）：} 当 $\lambda=0$ 或 $\vec{a}=\vec{0}$ 时，要单独考虑方向是否有定义，不能直接说同向或反向。
	\end{itemize}
\end{summarybox}